% Part: model-theory
% Chapter: models-of-arithmetic
% Section: models-of-q

\documentclass[../../../include/open-logic-section]{subfiles}

\begin{document}

\olfileid[fr]{mod}{mar}{mdq}
\section{Modèles de $\Th{Q}$}

\begin{explain}
Nous savons qu'il existe des !!{structure}s non standards
qui rendent vrais les mêmes !!{sentence}s que~$\Struct{N}$,
c'est-à-dire des modèles de~$\Th{TA}$. Puisque
$\Sat{N}{\Th{Q}}$, tout modèle de~$\Th{TA}$ est
aussi un modèle de~$\Th{Q}$. Cette dernière théorie
est beaucoup plus faible : par exemple,
$\Th{Q}\Proves/
\lforall[x][\lforall[y][\eq[(x+y)][(y+x)]]]$.
Une théorie plus faible a davantage de modèles.
Certains modèles de~$\Th{Q}$ rendent fausse la
commutativité de l'addition et ne peuvent donc
être des modèles de~$\Th{TA}$, ni d'ailleurs
de~$\Th{PA}$. Les modèles de~$\Th{Q}$ sont
aussi assez simples pour que nous puissions
en donner des descriptions explicites.
\end{explain}

\begin{ex}
\ollabel{ex:model-K-of-Q}
Considérons la !!{structure}~$\Struct{K}$ de domaine
$\Domain{K}=\Nat\cup\{a\}$, avec les interprétations
suivantes :
\begin{align*}
  \Assign{\Obj{0}}{K} & = 0\\
  \Assign{\prime}{K}(x) & =
  \begin{cases}
    x+1 & \text{si $x\in\Nat$}\\
    a & \text{si $x=a$}
  \end{cases}\\
  \Assign{+}{K}(x,y) & =
  \begin{cases}
    x+y & \text{si $x,y\in\Nat$}\\
    a & \text{sinon}
  \end{cases}\\
  \Assign{\times}{K}(x,y) & =
  \begin{cases}
    xy & \text{si $x,y\in\Nat$}\\
    0 & \text{si $x=0$ ou $y=0$}\\
    a & \text{sinon}
  \end{cases}\\
  \Assign{<}{K} & =
  \Setabs{\tuple{x,y}}{x,y\in\Nat\text{ et }x<y}
  \cup\Setabs{\tuple{x,a}}{x\in\Domain{K}}.
\end{align*}
Pour montrer que $\Sat{K}{\Th{Q}}$, il faut
vérifier que tous les axiomes de~$\Th{Q}$ sont
vrais dans~$\Struct{K}$. Écrivons
$x^\nssucc$ pour $\Assign{\prime}{K}(x)$,
le « successeur » de~$x$ dans~$\Struct{K}$;
$x\nsplus y$ pour $\Assign{+}{K}(x,y)$,
sa « somme »; $x\nstimes y$ pour
$\Assign{\times}{K}(x,y)$, son « produit »;
et $x\nsless y$ pour
$\tuple{x,y}\in\Assign{<}{K}$.
Ces abréviations permettent de présenter
plus clairement les opérations :
\[
\begin{array}{c|c}
  x & x^\nssucc \\
  \hline
  n & n+1 \\
  a & a
\end{array}
\qquad
\begin{array}{c|ccc}
  x\nsplus y & 0 & m & a \\
  \hline
  0 & 0 & m & a \\
  n & n & n+m & a \\
  a & a & a & a
\end{array}
\qquad
\begin{array}{c|ccc}
  x\nstimes y & 0 & m & a \\
  \hline
  0 & 0 & 0 & 0 \\
  n & 0 & nm & a \\
  a & 0 & a & a
\end{array}
\]
Dans les tableaux de l'addition et de la multiplication ci-dessus, $n,m$ désignent des entiers strictement positifs.\footnote{Précision éditoriale : les lignes et colonnes $0$ sont séparées ; sans cette convention, la case $a\nstimes m=a$ serait fausse pour $m=0$.} Pour $n,m\in\Nat$, on a $n\nsless m$
si et seulement si $n<m$; de plus,
$x\nsless a$ pour tout $x\in\Domain{K}$.

L'axiome
$\Sat{K}{\lforall[x][\lforall[y][
(\eq[x'][y']\lif\eq[x][y])]]}$
est satisfait, car $\nssucc$ est !!{injective}.
De même,
$\Sat{K}{\lforall[x][\eq/[\Obj 0][x']]}$,
car $0$ n'est le $\nssucc$-successeur d'aucun
élément de~$\Struct{K}$. Enfin,
$\Sat{K}{\lforall[x][(\eq[x][\Obj 0]\lor
\lexists[y][\eq[x][y']])]}$ :
pour $n>0$, $n=(n-1)^\nssucc$,
et $a=a^\nssucc$.

On a aussi
$\Sat{K}{\lforall[x][\eq[(x+\Obj 0)][x]]}$,
puisque $n\nsplus 0=n+0=n$ et
$a\nsplus 0=a$. Pour l'axiome
$\Sat{K}{\lforall[x][\lforall[y][
\eq[(x+y')][(x+y)']]]}$, distinguons les cas.
Si $n,m$ sont standards, alors
\begin{align*}
(n\nsplus m^\nssucc)
  &=(n+(m+1))=(n+m)+1
   =(n\nsplus m)^\nssucc\\
\intertext{En effet, $\nsplus$ et $^\nssucc$
  coïncident avec $+$ et $\prime$ sur les
  nombres standards. Pour tout
  $x\in\Domain{K}$, on a aussi}
(x\nsplus a^\nssucc)
  &=(x\nsplus a)=a=a^\nssucc
   =(x\nsplus a)^\nssucc\\
\intertext{Il reste le cas $x=a$ et $y$ standard,
  ainsi que le cas $x=y=a$ :}
(a\nsplus n^\nssucc)
  &=(a\nsplus(n+1))=a=a^\nssucc
   =(a\nsplus n)^\nssucc\\
(a\nsplus a^\nssucc)
  &=(a\nsplus a)=a=a^\nssucc
   =(a\nsplus a)^\nssucc.
\end{align*}
Ces distinctions sont même plus détaillées
que nécessaire : comme $a\nsplus z=a$
pour tout~$z$, on peut conclure directement
que $a\nsplus a^\nssucc=a$. On vérifie
les autres axiomes de la même manière.\footnote{Correction éditoriale : dans la discussion du modèle $\Struct{K}$, la source écrit $y=b$ alors que le seul élément ajouté est $a$. Le cas non standard est donc $y=a$.}

Ainsi, $\Struct{K}$ est un modèle de~$\Th{Q}$.
Son « addition » $\nsplus$ est même commutative.
Il existe pourtant des !!{sentence}s vrais
dans~$\Struct{N}$ et faux dans~$\Struct{K}$,
et inversement. En particulier,
$a\nsless a$, donc
$\Sat{K}{\lexists[x][x<x]}$ et
$\Sat/{K}{\lforall[x][\lnot x<x]}$.
Il s'ensuit que
$\Th{Q}\Proves/\lforall[x][\lnot x<x]$.
\end{ex}

\begin{prob}
Montrer que la !!{structure}~$\Struct{K}$ de
l'\olref[mod][mar][mdq]{ex:model-K-of-Q}
satisfait les axiomes restants de~$\Th{Q}$ :
\begin{align*}
  & \lforall[x][\eq[(x\times\Obj 0)][\Obj 0]]
    \tag{$!Q_6$}\\
  & \lforall[x][\lforall[y][
    \eq[(x\times y')][((x\times y)+x)]]]
    \tag{$!Q_7$}\\
  & \lforall[x][\lforall[y][
    (x<y\liff\lexists[z][\eq[(z'+x)][y]])]]
    \tag{$!Q_8$}.
\end{align*}
Trouver un !!{sentence} ne faisant intervenir,
parmi les symboles non logiques, que~$\prime$,
qui soit vrai dans~$\Struct{N}$ mais faux
dans~$\Struct{K}$.
\end{prob}

\begin{ex}
\ollabel{ex:model-L-of-Q}
Considérons maintenant la !!{structure}~$\Struct{L}$
de domaine $\Domain{L}=\Nat\cup\{a,b\}$.
Les interprétations
$\Assign{\prime}{L}=\nssucc$ et
$\Assign{+}{L}=\nsplus$ sont données par
\[
\begin{array}{c|c}
  x & x^\nssucc \\
  \hline
  n & n+1 \\
  a & a\\
  b & b
\end{array}
\qquad
\begin{array}{c|ccc}
  x\nsplus y & m & a & b\\
  \hline
  n & n+m & b & a\\
  a & a & b & a\\
  b & b & b & a
\end{array}
\]
La fonction $\nssucc$ est !!{injective},
$0$ n'est pas dans son image, et tous les
éléments de~$\Domain{L}$ autres que~$0$
y figurent. Les axiomes $!Q_1$--$!Q_3$
sont donc vrais dans~$\Struct{L}$.
Comme $x\nsplus 0=x$ pour tout~$x$,
$!Q_4$ l'est également.

Pour $!Q_5$, comparons
$x\nsplus y^\nssucc$ et
$(x\nsplus y)^\nssucc$.
Ils sont égaux si $x$ et $y$ sont
standards, car les opérations y sont
les opérations usuelles. Si $x$ est
non standard et $y$ standard,
$x\nsplus y^\nssucc=x=x^\nssucc
=(x\nsplus y)^\nssucc$.
Si les deux sont non standards,
il y a quatre cas :
\begin{align*}
& a\nsplus a^\nssucc
   =b=b^\nssucc=(a\nsplus a)^\nssucc\\
& b\nsplus b^\nssucc
   =a=a^\nssucc=(b\nsplus b)^\nssucc\\
& b\nsplus a^\nssucc
   =b=b^\nssucc=(b\nsplus a)^\nssucc\\
& a\nsplus b^\nssucc
   =a=a^\nssucc=(a\nsplus b)^\nssucc\\
\intertext{Enfin, si $x=n$ est standard et
  $y$ non standard,}
& n\nsplus a^\nssucc
   =n\nsplus a=b=b^\nssucc
   =(n\nsplus a)^\nssucc\\
& n\nsplus b^\nssucc
   =n\nsplus b=a=a^\nssucc
   =(n\nsplus b)^\nssucc.
\end{align*}
Ainsi, $\Sat{L}{!Q_5}$.\footnote{Correction éditoriale : dans le cas $b\nsplus a^\nssucc$, la dernière expression contient $a$, et non la variable $y$ laissée par erreur dans la source.} Cependant,
$a\nsplus 0\neq 0\nsplus a$, donc
$\Sat/{L}{\lforall[x][\lforall[y][
\eq[(x+y)][(y+x)]]]}$.
\end{ex}

\begin{prob}
Compléter la !!{structure}~$\Struct{L}$ de
l'\olref[mod][mar][mdq]{ex:model-L-of-Q}
par des interprétations $\nstimes$ et
$\nsless$ de~$\times$ et de~$<$.
Montrer que la !!{structure} obtenue
satisfait les axiomes restants de~$\Th{Q}$ :
\begin{align*}
& \lforall[x][\eq[(x\times\Obj 0)][\Obj 0]]
  \tag{$!Q_6$}\\
& \lforall[x][\lforall[y][
  \eq[(x\times y')][((x\times y)+x)]]]
  \tag{$!Q_7$}\\
& \lforall[x][\lforall[y][
  (x<y\liff\lexists[z][\eq[(z'+x)][y]])]]
  \tag{$!Q_8$}.
\end{align*}
\end{prob}

\begin{prob}
Dans $\Struct{L}$ de
l'\olref[mod][mar][mdq]{ex:model-L-of-Q},
$a^\nssucc=a$ et $b^\nssucc=b$.
Existe-t-il un modèle de~$\Th{Q}$
où $a^\nssucc=b$ et
$b^\nssucc=a$ ?
\end{prob}

\begin{explain}
Nous avons construit explicitement des
modèles de~$\Th{Q}$ où les !!{element}s
non standards se trouvent « au-delà »
des éléments standards. Les axiomes
l'imposent : un !!{element} non standard~$x$
ne peut satisfaire $x\nsless 0$,
car $\Th{Q}\Proves
\lforall[x][\lnot x<0]$
(voir le \olref[inc][req][min]{lem:less-zero}).
De plus, pour tout~$n$,
\[
\Th{Q}\Proves
\lforall[x][(x<\num{n}'\lif
(\eq[x][\num{0}]\lor\eq[x][\num{1}]
\lor\dots\lor\eq[x][\num{n}]))]
\]
(voir le \olref[inc][req][min]{lem:less-nsucc});
aucun élément non standard ne peut donc
être $\nsless$-inférieur à un numéral
standard positif.
\end{explain}

\end{document}
